\chapter{Results}

The results of the research are presented in this chapter. The first section presents the results of the corpus-based metrics, while the second the results of the morph-based metrics. 

\section{Corpus based metrics}

Three metrics are the most important in this category:
\begin{itemize}
    \item type to token ratio ($TTR$),
    \item word-type entropy ($H$),
    \item information in word structure ($IiS$).
\end{itemize}

\subsection{Quality of the result}

\subsubsection{Comparison with previous work}

To evaluate whether the results are consistent with previous work, they were compared with the work of \textcite{Coltekin2022}, using TTR and entropy. The Spearman rank correlation was computed over the languages shared between the original paper and this thesis.
\begin{itemize}
    \item Type to token ratio ($\text{TTR}$). The Spearman rank correlation is $0.630$. The languages whose rank changed the most are Uyghur ($32$\footnote{The rank in \textcite{Coltekin2022} minus the rank in this thesis. A positive value means that the metric is higher in this thesis.}), Latin ($26$) and Latvian ($-18$).
    \item Entropy ($H$). The Spearman rank correlation is $0.734$. The languages whose rank changed the most are Hebrew ($34$), Latin ($24$) and Uyghur ($17$).
    \item Information in word structure (IiS). The Spearman rank correlation is $0.313$. The languages whose rank changed the most are Chinese ($-36$), Indonesian ($35$) and Vietnamese ($28$).
\end{itemize}

The correlation for entropy is very high, suggesting that the results are very similar. Even the TTR correlation is still higher than the typical correlation among the different complexity measures in that paper. Latin and Uyghur are ranked much higher on both metrics in this thesis. The reason for the difference in Latin may be the variety of sources for this dead language and the lack of a single standard for what Latin is. The difference in Uyghur may be that \textcite{Coltekin2022} use the Latin script for it, whereas in this thesis Uyghur is in Arabic script. Overall, the results for these two metrics are very consistent.

Information in word structure has the lowest correlation. This can be explained by the fact that the methodology used here is completely different. The text is generated by random sampling and character sizes are unified across scripts.

\subsubsection{Related languages}
Another way of evaluating whether the metrics are relevant is to compare languages that should have similar complexity based on their genealogical proximity. The most closely related languages are presented in the following two tables. Table~\ref{tab:duplicate_languages} lists dialects and variants of the same language. Table~\ref{tab:closely_related_languages} lists groups of closely related and mutually intelligible languages. These languages are expected to lie close to each other in the graphs.
    
\begin{table}[p]
\centering
\footnotesize
\caption{Groups of languages treated as duplicates or near-duplicates}
\label{tab:duplicate_languages}
\begin{tabularx}{\textwidth}{
>{\hsize=0.6\hsize\raggedright\arraybackslash}X|
>{\hsize=1.4\hsize\raggedright\arraybackslash}X
}
\textbf{Group} & \textbf{Reason} \\
\hline
\hline
Norwegian (nno, nob, nor) & Nynorsk, Bokmål, and generic Norwegian -- orthographic variants of Norwegian \\
\hline
Serbo-Croatian (hbs, hrv, bos, srp) & One language, four political standards \\
\hline
Malay--Indonesian (ind, msa, zsm) & Indonesian is standardized Malay; msa and zsm are macro/specific codes \\
\hline
Hindi--Urdu (hin, urd) & Same grammar, different script \\
\hline
Persian (fas, pes) & fas = macrolanguage, pes = Iranian Persian \\
\hline
Estonian (est, ekk) & est = macrolanguage, ekk = Standard Estonian \\
\hline
Latvian (lav, lvs) & lav = macrolanguage, lvs = Standard Latvian \\
\hline
Azerbaijani (aze, azj) & aze = macrolanguage, azj = North Azerbaijani \\
\hline
Uzbek (uzb, uzn) & uzb = macrolanguage, uzn = Northern Uzbek \\
\hline
Malagasy (mlg, plt) & mlg = macrolanguage, plt = Plateau Malagasy \\
\hline
Swahili (swa, swh) & swa = macrolanguage, swh = specific Swahili \\
\hline
Oriya (ori, ory) & Same language, duplicate ISO codes \\
\hline
Tagalog--Filipino (tgl, fil) & Filipino is standardized Tagalog \\
\hline
Albanian (als, sqi) & als = Tosk Albanian, sqi = macrolanguage \\
\hline
Punjabi (pan, pnb) & Eastern and Western Punjabi; very closely related \\
\hline
\end{tabularx}

 \vspace{1em}
\end{table}
\begin{table}[p]
\caption{Groups of very closely related languages expected to correlate strongly}
\label{tab:closely_related_languages}
\begin{tabularx}{\textwidth}{
>{\hsize=0.8\hsize\raggedright\arraybackslash}X|
>{\hsize=1.2\hsize\raggedright\arraybackslash}X
}
\textbf{Group} & \textbf{Reason} \\
\hline
\hline
Czech--Slovak (ces, slk) & Mutually intelligible West Slavic languages \\
\hline
Bulgarian--Macedonian (bul, mkd) & Very closely related South Slavic languages with highly similar grammatical systems \\
\hline
Danish--Swedish--Norwegian (dan, swe, nno/nob) & Mainland Scandinavian languages with mutual intelligibility \\
\hline
East Slavic (rus, ukr, bel) & Closely related East Slavic languages with similar morphological systems \\
\hline
Spanish--Portuguese--Galician (spa, por, glg) & Closely related Romance languages; Galician and Portuguese are especially close \\
\hline
Turkish--Azerbaijani--Turkmen (tur, aze/azj, tuk) & Oghuz Turkic languages with similar agglutinative morphology \\
\hline
Kazakh--Kyrgyz--Tatar--Bashkir (kaz, kir, tat, bak) & Kipchak Turkic languages with related morphological systems \\
\hline
\end{tabularx}
\end{table}

Figure~\ref{fig:ttrToEntropy} shows these groups in color, connected by lines. The $x$ axis shows TTR on a logarithmic scale, the $y$ axis shows entropy, and the size of each point represents the original token count of the dataset. There is clearly a strong relationship between the two metrics, with a correlation of $0.89$.

It is visible at first glance that the highlighted languages do indeed cluster together. For most languages, the distance on the $y$ axis is smaller than the distance on the $x$ axis, which suggests that entropy is the better metric. It is also a more sophisticated metric and shows a higher correlation with \textcite{Coltekin2022}.

The most positive observation is the very low correlation with the original token count for all the metrics. It is below $0.08$ in absolute value for each of them, which was the main goal of preprocessing.

It is also worth noting that there is almost no correlation between entropy and information in word structure.

Figure~\ref{fig:entropyToCompressibility} is very similar, but on the $x$ axis it plots information in word structure (compressibility). The following observations can be made.
\begin{itemize}
    \item Some languages are closer to each other on the $x$ axis and others on the $y$ axis, so both metrics carry useful information.
    \item Languages with different scripts are marked in red. The script clearly has a substantial influence on the information in word structure. This is visible for the pair Hindi and Urdu, and for the pair Punjabi and Western Punjabi. The variant written in Arabic script is always less compressible.
    \item On the other hand, Uzbek and Northern Uzbek look comparable even though they are written in different scripts, which suggests that Cyrillic and Latin may be more comparable on this metric. Unfortunately, there is only this one example.
    \item There is a clear error in the Malagasy group, since the two languages are very far from each other on both metrics.
\end{itemize}


\begin{figure}[p]
    \centering
        \includegraphics[width=1\linewidth]{img/auto/ttrToEntropy.pdf}
        \caption{TTR to Entropy with groups of very similar languages (sizes of languages representing size of original data)}
        \label{fig:ttrToEntropy}
\end{figure}


\begin{figure}[p]
    \centering

    \includegraphics[width=1\linewidth]{img/auto/entropyToCompressibility.pdf}
    \caption{Information in word structure to Entropy with groups of very similar languages (sizes of languages representing size of original data).}
    \label{fig:entropyToCompressibility}
\end{figure}

\newpage
\subsubsection{Indo-European languages}

It may also be interesting to look at the largest family present in the data. The Indo-European languages are shown in Figure~\ref{fig:ie_entropy_compressibility}, in order to see whether the metrics produce the expected family clustering. Clear clustering of the families can be seen.

Most of the \textbf{Slavic} languages are in the high-entropy part of the graph, probably because of their rich inflection.
\begin{itemize}
    \item East Slavic languages are very close to each other, and they sit between the related Baltic languages.
    \item West Slavic languages are very close in entropy, but Czech appears to be less compressible. Upper Sorbian, on the other hand, has slightly lower entropy, which could be caused by the smaller number of speakers and written materials.
    \item South Slavic languages cluster by script. Bosnian, Serbo-Croatian and Slovenian are closer together, while Serbian, Bulgarian and Macedonian, written in Cyrillic, sit further away from them.
\end{itemize}

\textbf{Germanic} languages form a recognisable cluster in the middle of the graph.
\begin{itemize}
    \item English is curiously placed in the middle of them, although it differs significantly from the rest of the Germanic languages.
    \item The most visible outliers are the German dialects Swiss German and Bavarian, which may indicate either lower data quality or that these unofficial dialects have unique features and lack a fully unified grammar.
    \item Afrikaans has one of the lowest entropies among all Indo-European languages, which is expected since it was created as a simplified version of Dutch.
    \item The opposite case is Luxembourgish, which sits at the top of the entropy scale.
\end{itemize}

\textbf{Romance} languages are interestingly arranged in a line. The higher the entropy, the lower the information in word structure, starting with Corsican, then the largest Romance languages, and ending with Walloon, which has the lowest entropy. One outlier is Latin, which has higher entropy. This is expected, since Latin has the richest inflection of all Romance languages. Latin also ranks higher in entropy here than in \textcite{Coltekin2022}. The most spoken Romance languages such as Italian, French and Spanish are clustered close together.

The \textbf{Indic} group has two entropy clusters. The first one, around Hindi, Urdu and Punjabi, has low entropy. These languages are close geographically and linguistically. The other group has high entropy and contains Bengali, Marathi, Nepali and Gujarati. These languages are also high on information in word structure.

It is also apparent from the graph that the script may influence not only the compressibility but also the entropy. For example, all languages written in Arabic script have lower entropy and lower compressibility than the rest of their group. As explained in Section~\ref{sec:scripts}, the reason may be that abjad scripts use the same orthography for all words and word-forms built from the same consonants. Figure~\ref{fig:compressibility_by_script_violin} clearly shows that languages cluster by script. This measure should therefore only be used to compare languages within the same script.

\begin{figure}[p]
    \centering
    \includegraphics[width=1\linewidth]{img/auto/ie_entropy_compressibility.pdf}
    \caption{Information in word structure to Entropy with groups only for indo-european languages. }
    \label{fig:ie_entropy_compressibility}
\end{figure}
\begin{figure}[p]
    \centering
    \includegraphics[width=1\linewidth]{img/auto/compressibility_by_script_violin.pdf}
    \caption{Languages according to information in word structure on $x$ axes. The languages are grouped on y axes by script type.}
    \label{fig:compressibility_by_script_violin}
\end{figure}

\subsubsection{Typology}

Using the typology presented in Chapter~\ref{chap:typology}, every language was assigned to one of the categories. Figure~\ref{fig:entropy_by_typology_violin} shows the distribution of entropy by typology. The agglutinative languages are clearly the most spread across the entropy axis, containing both languages with the highest and the lowest entropy. This may be caused by the difficulty of clearly distinguishing between categories, which leaves the agglutinative category as a fallback for borderline languages. It may also be caused by the fact that the definition of agglutinative does not specify how many morphs a word should contain on average, while the other typologies are more strict in this respect.

Interestingly, the fusional languages form two clusters. The first one is closer in entropy to the isolating languages, while the second one is closer to the largest cluster of agglutinative and introflexive languages.
%todo what are what

Apart from the agglutinative languages, the typology follows the expected order of complexity.

\begin{figure}[H]
    \centering
    \includegraphics[width=1\linewidth]{img/auto/entropy_by_typology_violin.pdf}
    \caption{Entropy by typology}
    \label{fig:entropy_by_typology_violin}
\end{figure}

\subsubsection{Conclusion}

\textbf{Entropy} appears to be the most robust metric. It places related languages close to each other and is not influenced by external factors. It is also the most consistent with the previous work of \textcite{Coltekin2022}.

\textbf{TTR} appears to be redundant with entropy. The two are highly correlated, but TTR generally performs worse, both with respect to the work of \textcite{Coltekin2022} and in clustering related languages.

\textbf{Information in word structure} is influenced by the script and is therefore more useful for comparing languages written in the same script.

\subsection{Zipf's slope}

Zipf's slope is the slope of the regression line in the log-log plot of rank against frequency. The theory says that this distribution of words is similar across languages and that the slope is close to $-1$. As Figure~\ref{fig:zipfSlope} shows, the slope is in fact between $-1.1$ and $-1.4$ for most languages. The outliers are Volapük, an artificial language, and Maltese, which does not have a high-quality dataset.

On the other end of the scale, the smallest slope belongs to Inuktitut. Inuktitut is a highly polysynthetic language. It has one of the smallest token counts but, even so, one of the highest type counts, which gives it by far the highest TTR and the 4\textsuperscript{th} highest entropy.


\begin{figure}[H]
    \centering
    \includegraphics[width=0.75\linewidth]{img/auto/zipfSlopeComparison.pdf}
    \caption{Zipf's slope representation with referenced Entropy}
    \label{fig:zipfSlope}
\end{figure}
The graph represents all the slopes starting from the same point $(0,0)$ and ending at $x=1$ and $y=\text{slope}$. On the right side there is the graph representing entropy for each language at the same $y$ as at the left graph.

\newpage
\section{Morph based metrics}

From the previous section it is clear that word entropy can be trusted to a reasonable degree. The first sanity check for the morph based metrics is therefore to compare them with word entropy. Figure~\ref{fig:morph_word_entropy} shows the relationship between the two. Colours represent language families. Two trends are visible in the graph.

Most of the Niger-Congo languages are low on both word and morph entropy. The exceptions are Zulu and Xhosa, which score high on word entropy but only slightly higher on morph entropy. Even so, there is a negative correlation between morph and word entropy in this group.

There is a clear issue with the MorphyNet datasets. Languages where MorphyNet is the only segmenter score very high on morph entropy in comparison with related languages. The most obvious example is Polish, where words such as \ttt{samowystarczalności} are routinely left unsegmented. The same applies to Serbo-Croatian. The same is true for Portuguese, whose morph entropy is much higher than that of Spanish. Compare the correctly segmented Spanish \ttt{ex+@clu+ir}\footnote{The \ttt{@} symbol marks the morph identified as a root.} with the unsegmented Portuguese \ttt{excluir}. MorphyNet datasets are the only ones that score below $1.5$ morphs per word, and they therefore have to be excluded from the analysis.

German and Russian both have two segmenters in the datasets that produce reasonable results. In both cases, the more conservative segmentation leads to higher morph entropy. The difference is around $0.5$ points across the board, which can be used as a rough estimate of the error even for the better segmenters.

\begin{figure}[H]
    \centering
    \includegraphics[width=1\linewidth]{img/auto/morph_word_entropy.pdf}
    \caption{Morph Entropy vs Word Entropy}
    \label{fig:morph_word_entropy}
\end{figure}

\subsection{Morph typology metrics}

How data was segmented can be visible in Figure~\ref{fig:morphs}. The languages are sorted by the average number of morphs per word type. The yellow bar represents the average number of suffixes. The blue and red bars represent the other affixes. The green part represents the average number of roots minus one, since there is always at least one root, so this is in fact the additional root count. Slavic languages score high, as do Kannada, Georgian and Finnish, all known for their rich morphology. Languages whose morphology comes from MorphyNet are predictably very low. The limitations of the morph typology are also clearly visible. For example, Finnish is known for not having prefixes, and Belarusian is reported to have two roots on average, which does not match reality.

\begin{figure}[hbtp]
    \centering
    \includegraphics[width=1\linewidth]{img/auto/morphs.pdf}
    \caption{Average numbers of prefixes, roots, interfixes, and suffixes per word across languages, with point markers indicating script and language family and the right panel showing morphological typology}
    \label{fig:morphs}
\end{figure}

A similar graph could be produced from the average morph-type count per token. The ordering of the languages would be very similar, but this approach also favours morphologically less interesting words (which tend to be shorter and more frequent) over the complex words, which usually have lower frequencies.

The absolute values are not correct, but the data are still comparable. Isolating

languages have the lowest scores, less than one suffix on average and a negligible amount of other morphs.

Figure~\ref{fig:compounding} shows the compounding score. All languages with a compounding score higher than $1$, that is, those preferring compounding over affixation, are agglutinative languages (with the exception of Armenian, which is closest to $1$). However, the languages with the lowest compounding score, close to zero, are also agglutinative, mainly Niger-Congo and Austronesian languages. Most of the fusional and isolating languages sit in the middle, but they prefer affixation.

\begin{figure}[hbtp]
    \centering
    \includegraphics[width=1\linewidth]{img/auto/compounding.pdf}
    \caption{Average compounding index across languages, with point markers indicating script and language family and the right panel showing morphological typology.}
    \label{fig:compounding}
\end{figure}

Agglutinative languages should have more morphs classified as roots. This is true for most of them, with the exception of the low-entropy Niger-Congo languages, Tagalog and the Altaic languages.

Fusional languages should have more morphs classified as affixes.



\subsubsection{Affix deviation}

According to the affix deviation results, most languages prefer suffixation over prefixation. Slavic languages appear to be the most suffix-leaning, which can be explained by their large number of suffixes. On the other hand, Tagalog, German and Turkish appear to prefer prefixation, which is unlikely to be true. In German, this can be caused by the annotator treating some frequent compounding roots as prefixes, as in \ttt{kauf+ent+@scheid+ung}, but mainly by treating endings such as \ttt{e} as roots. Tagalog, on the other hand, does in fact use many prefixes. As before, the absolute values are not perfect, but the data are still comparable.

\subsection{Conclusion}

For the following comparison, a grouping metric was designed.

\begin{itemize}
    \item Standardise all values for the metric to remove the influence of different distributions.
    \item Normalise all values for the metric so they lie in the range $0$ to $1$.
    \item Compute the squared average distance of two languages within one group.
    \item Average across all groups.
    \item Apply $1-n$ to flip the direction so that a higher value means better grouping.
    \item Since all values are above $0.8$, they were stretched by the formula $\frac{n-0.8}{0.2}$ to make the differences more visible.
\end{itemize}

With this formula, the grouping score for the segmentation metrics was calculated for the following groups.
\begin{itemize}
    \item Family and genus.
    \item Script grouping. This is a negative control, since the script is not a linguistic feature.
    \item Morphological typology groups.
\end{itemize}

Figure~\ref{fig:morphology_grouping} shows the grouping scores together with the correlation with word entropy and with the original token count, the latter being another negative control.


\begin{figure}[hbtp]
    \centering
    \includegraphics[width=1\linewidth]{img/auto/morphology_grouping.pdf}
    \caption{Grouping-property comparison matrix for complexity metrics and correlation with word entropy and original token count}
    \label{fig:morphology_grouping}
\end{figure}

The first important observation is that none of the metrics correlate with the original token count, which was the main goal of preprocessing.

Three metrics group related languages better than word entropy does. These are information in word structure, affix deviation and affix entropy. The price for this, however, is a much higher dependence on script. Morph entropy on tokens looks more reliable than morph entropy on word types, which means that the morphs are weighted by the frequency of the word types. It is better at predicting typology and much better at predicting family grouping, with only a moderate increase in script dependence. It also has much higher correlation with word entropy.

The compounding index appears to be the best metric among the morphological typology metrics. It balances the dependence on script and on the original token count better than the other metrics.

Average morph count per word type is also a reasonable metric, as shown in the previous section.

In conclusion, the following metrics seem to be the most relevant.

\begin{itemize}
    \item Morph entropy on tokens.
    \item Average morph count per word type.
    \item Compounding index.
\end{itemize}

The results are far from perfect. It is hard to define what a perfect segmentation would mean, but the models have currently some obvious errors. For deeper analysis, more segmentation data would be needed, ideally from a similar source of training data. The data are strongly influenced by the quality of segmentation and by the fact that very few languages are segmented at all.

% Coclusion